% ============================================================
%  PAQUETES
%  El orden importa: babel y amsmath van ANTES de hyperref, y
%  hyperref va de último (parchea comandos de los demás).
% ============================================================
\usepackage[utf8]{inputenc}
\usepackage[T1]{fontenc}
\usepackage{lmodern}

% Español: silabación correcta y nombres traducidos ("Índice").
%   es-nolayout     -> no cambia sangrías ni pone punto tras el número de sección
%   es-nolists      -> no cambia las viñetas de itemize/enumerate
%   es-noshorthands -> ~ y " NO se vuelven caracteres activos (romperían listings)
\usepackage[spanish,es-nolayout,es-nolists,es-noshorthands]{babel}
\unaccentedoperators   % \bmod, \lim, \max... se escriben "mod", "lim" y no "mód", "lím"

\usepackage{xcolor}
\usepackage{graphicx}
\usepackage{amsmath}
\usepackage{amssymb}
\usepackage{listings}
\usepackage{titlesec}
\usepackage{fancyhdr}
\usepackage[hidelinks]{hyperref}   % hidelinks: enlaces clicables sin cajas de color

\graphicspath{{images/}}

% ------------------------------------------------------------
%  Símbolos Unicode permitidos en el TEXTO corrido.
%  Se definen con \ensuremath para que también funcionen si
%  quedan dentro de $...$ o de \bigO{...}. NO usarlos dentro
%  de lstlisting (listings no los entiende).
% ------------------------------------------------------------
\DeclareUnicodeCharacter{2264}{\ensuremath{\leq}}            % ≤
\DeclareUnicodeCharacter{2265}{\ensuremath{\geq}}            % ≥
\DeclareUnicodeCharacter{2260}{\ensuremath{\neq}}            % ≠
\DeclareUnicodeCharacter{2261}{\ensuremath{\equiv}}          % ≡
\DeclareUnicodeCharacter{2192}{\ensuremath{\rightarrow}}     % →
\DeclareUnicodeCharacter{2212}{\ensuremath{-}}               % − (signo menos)
\DeclareUnicodeCharacter{00D7}{\ensuremath{\times}}          % ×
\DeclareUnicodeCharacter{00B1}{\ensuremath{\pm}}             % ±
\DeclareUnicodeCharacter{00B2}{\ensuremath{^{2}}}            % ²
\DeclareUnicodeCharacter{00B7}{\ifmmode\cdot\else\textperiodcentered\fi} % ·
\DeclareUnicodeCharacter{03A3}{\ensuremath{\Sigma}}          % Σ
\DeclareUnicodeCharacter{00AA}{\textordfeminine}             % ª
\DeclareUnicodeCharacter{0151}{\H{o}}                        % ő (Kőnig)
\DeclareUnicodeCharacter{2013}{--}                           % –
\DeclareUnicodeCharacter{2014}{---}                          % —
\DeclareUnicodeCharacter{2026}{\ldots}                       % …
